
\chapter{2017年北京卷（文）}

\section{单选题}

\begin{problem}
已知全集 \(U=\R\)，集合 \(A=\{x\mid x<-2\text{或}x>2\}\)，则 \(C_UA=\)（\quad）.

\choices
    {\(\left(-2,2\right)\)}
    {\(\left(-\infty,-2\right)\cup\left(2,+\infty\right)\)}
    {\(\left[-2,2\right]\)}
    {\(\left(-\infty,-2\right]\cup\left[2,+\infty\right)\)}
\end{problem}

\begin{answer}
C
\end{answer}

\begin{solution}
由集合 \(A\) 的定义，\(A=(-\infty,-2)\cup(2,+\infty)\). 因为全集为 \(\R\)，所以 \(C_UA=U\mathbin{\backslash}A=[-2,2]\). 故选 C.
\end{solution}

\begin{problem}
若复数 \((1-\i)(a+\i)\) 在复平面内对应的点在第二象限，则实数 \(a\) 的取值范围是（\quad）.

\choices
    {\(\left(-\infty,1\right)\)}
    {\(\left(-\infty,-1\right)\)}
    {\(\left(1,+\infty\right)\)}
    {\(\left(-1,+\infty\right)\)}
\end{problem}

\begin{answer}
B
\end{answer}

\begin{solution}
因为 \(\i^2=-1\)，所以
\((1-\i)(a+\i)=a+\i-a\i-\i^2=(a+1)+(1-a)\i\).
该复数对应的点在第二象限，当且仅当实部和虚部分别满足 \(a+1<0\) 与 \(1-a>0\)，即 \(a<-1\) 且 \(a<1\). 因此 \(a<-1\)，故选 B.
\end{solution}

\begin{problem}
执行如图所示的程序框图，输出的 \(S\) 值为（\quad）

\texfigure[max width=6.0cm]{img_repaint/2017/beijing_liberal/q3_fig1.tex}

\choices
    {\(2\)}
    {\(\dfrac32\)}
    {\(\dfrac53\)}
    {\(\dfrac85\)}
\end{problem}

\begin{answer}
C
\end{answer}

\begin{solution}
初始时 \(k=0\)，\(s=1\). 因为 \(k<3\)，第一次循环后 \(k=1\)，\(s=\dfrac{1+1}{1}=2\)；第二次循环后 \(k=2\)，\(s=\dfrac{2+1}{2}=\dfrac32\)；第三次循环后 \(k=3\)，\(s=\dfrac{\frac32+1}{\frac32}=\dfrac53\).
此时 \(k<3\) 不成立，程序输出 \(s=\dfrac53\)，故选 C.
\end{solution}

\begin{problem}
若 \(x\)，\(y\) 满足 \(\begin{cases}
x\le3,\\
x+y\ge2,\\
y\le x,
\end{cases}\) 则 \(x+2y\) 的最大值为（\quad）.

\choices
    {\(1\)}
    {\(3\)}
    {\(5\)}
    {\(9\)}
\end{problem}

\begin{answer}
D
\end{answer}

\begin{solution}
由 \(x+y\ge2\) 和 \(y\le x\) 得 \(2\le x+y\le2x\)，所以 \(x\ge1\). 又因为 \(y\le x\)，故 \(x+2y\le3x\le9\)，其中最后一个不等号使用了 \(x\le3\).
当 \(x=3\)，\(y=3\) 时，三个条件均成立，且 \(x+2y=9\). 因此最大值为 \(9\)，故选 D.
\end{solution}

\begin{problem}
已知函数 \(f(x)=3^x-\left(\dfrac13\right)^x\)，则 \(f(x)\)（\quad）.

\choices
    {是偶函数，且在 \(\R\) 上是增函数}
    {是奇函数，且在 \(\R\) 上是增函数}
    {是偶函数，且在 \(\R\) 上是减函数}
    {是奇函数，且在 \(\R\) 上是减函数}
\end{problem}

\begin{answer}
B
\end{answer}

\begin{solution}
将 \(\left(\dfrac13\right)^x\) 写成 \(3^{-x}\)，则 \(f(x)=3^x-3^{-x}\). 于是 \(f(-x)=3^{-x}-3^x=-f(x)\)，所以 \(f\) 是奇函数.
又 \(f'(x)=\ln3\cdot3^x+\ln3\cdot3^{-x}>0\)，故 \(f(x)\) 在 \(\R\) 上是增函数. 因此选 B.
\end{solution}

\begin{problem}
某三棱锥的三视图如图所示，则该三棱锥的体积为（\quad）

\texfigure[max width=6.0cm]{img_repaint/2017/beijing_liberal/q6_fig1.tex}

\choices
    {\(60\)}
    {\(30\)}
    {\(20\)}
    {\(10\)}
\end{problem}

\begin{answer}
D
\end{answer}

\begin{solution}
由三视图可知，该三棱锥的底面为直角三角形，两条直角边长分别为 \(5\) 和 \(3\)，棱锥的高为 \(4\). 所以其体积为
\(V=\dfrac13\times\left(\dfrac12\times5\times3\right)\times4=10\).
故选 D.
\end{solution}

\begin{problem}
设 \(\bs{m}\)，\(\bs{n}\) 为非零向量，则“存在负数 \(\lambda\)，使得 \(\bs{m}=\lambda\bs{n}\)”是“\(\bs{m}\cdot\bs{n}<0\)”的（\quad）.

\choices
    {充分而不必要条件}
    {必要而不充分条件}
    {充分必要条件}
    {既不充分也不必要条件}
\end{problem}

\begin{answer}
A
\end{answer}

\begin{solution}
若存在负数 \(\lambda\) 使 \(\bs{m}=\lambda\bs{n}\)，则
\(\bs{m}\cdot\bs{n}=\lambda\bs{n}\cdot\bs{n}=\lambda\lvert\bs{n}\rvert^2<0\)，所以该条件是充分条件.
但 \(\bs{m}\cdot\bs{n}<0\) 不一定使两向量共线. 例如取 \(\bs{m}=(1,0)\)，\(\bs{n}=(-1,1)\)，则 \(\bs{m}\cdot\bs{n}=-1<0\)，而 \(\bs{m}\) 不是 \(\bs{n}\) 的实数倍. 因此该条件充分而不必要，故选 A.
\end{solution}

\begin{problem}
根据有关资料，围棋状态空间复杂度的上限 \(M\) 约为 \(3^{361}\)，而可观测宇宙中普通物质的原子总数 \(N\) 约为 \(10^{80}\). 则下列各数中与 \(\dfrac{M}{N}\) 最接近的是（\quad）.

（参考数据：\(\lg 3\approx0.48\)）

\choices
    {\(10^{33}\)}
    {\(10^{53}\)}
    {\(10^{73}\)}
    {\(10^{93}\)}
\end{problem}

\begin{answer}
D
\end{answer}

\begin{solution}
由 \(\lg3\approx0.48\)，得
\(\lg\dfrac{M}{N}=361\lg3-80\approx361\times0.48-80=93.28\).
所以 \(\dfrac{M}{N}\) 的数量级约为 \(10^{93}\)，在所给各数中与之最接近的是 \(10^{93}\)，故选 D.
\end{solution}

\section{填空题}

\begin{problem}
在平面直角坐标系 \(xOy\) 中，角 \(\alpha\) 与角 \(\beta\) 均以 \(Ox\) 为始边，它们的终边关于 \(y\) 轴对称. 若 \(\sin\alpha=\dfrac13\)，则 \(\sin\beta=\fillinblank{}\).
\end{problem}

\begin{answer}
\(\dfrac13\)
\end{answer}

\begin{solution}
终边关于 \(y\) 轴对称时，两条终边在单位圆上的纵坐标相同，因此正弦值相同. 所以 \(\sin\beta=\sin\alpha=\dfrac13\).
\end{solution}

\begin{problem}
若双曲线 \(x^2-\dfrac{y^2}{m}=1\) 的离心率为 \(\sqrt3\)，则实数 \(m=\fillinblank{}\).
\end{problem}

\begin{answer}
2
\end{answer}

\begin{solution}
双曲线可写为标准形式 \(\dfrac{x^2}{1}-\dfrac{y^2}{m}=1\)，故 \(a^2=1\)，\(b^2=m\)，并且 \(m>0\). 其焦半距满足 \(c^2=a^2+b^2=1+m\)，离心率为
\(e=\dfrac ca=\sqrt{1+m}=\sqrt3\).
于是 \(1+m=3\)，解得 \(m=2\).
\end{solution}

\begin{problem}
已知 \(x\ge0\)，\(y\ge0\)，且 \(x+y=1\)，则 \(x^2+y^2\) 的取值范围是\fillinblank{}.
\end{problem}

\begin{answer}
\(\left[\dfrac12,1\right]\)
\end{answer}

\begin{solution}
由 \(x+y=1\) 且 \(x\ge0\)，\(y\ge0\)，可令 \(y=1-x\)，其中 \(0\le x\le1\). 则
\(x^2+y^2=x^2+(1-x)^2=2\left(x-\dfrac12\right)^2+\dfrac12\).
当 \(0\le x\le1\) 时，\(\left(x-\dfrac12\right)^2\in\left[0,\dfrac14\right]\)，所以 \(x^2+y^2\in\left[\dfrac12,1\right]\).
\end{solution}

\begin{problem}
已知点 \(P\) 在圆 \(x^2+y^2=1\) 上，点 \(A\) 的坐标为 \((-2,0)\)，\(O\) 为原点，则 \(\overrightarrow{AO}\cdot\overrightarrow{AP}\) 的最大值为\fillinblank{}.
\end{problem}

\begin{answer}
6
\end{answer}

\begin{solution}
设 \(P=(x,y)\). 则 \(\overrightarrow{AO}=(2,0)\)，\(\overrightarrow{AP}=(x+2,y)\)，从而
\(\overrightarrow{AO}\cdot\overrightarrow{AP}=2(x+2)=2x+4\).
又 \(P\) 在单位圆上，所以 \(-1\le x\le1\). 因此 \(2x+4\le6\)，当 \(P=(1,0)\) 时取等号，故最大值为 \(6\).
\end{solution}

\begin{problem}
能够说明“设 \(a\)，\(b\)，\(c\) 是任意实数. 若 \(a>b>c\)，则 \(a+b>c\)”是假命题的一组整数 \(a\)，\(b\)，\(c\) 的值依次为\fillinblank{}.
\end{problem}

\begin{answer}
\(-1\)，\(-2\)，\(-3\)
\end{answer}

\begin{solution}
取 \(a=-1\)，\(b=-2\)，\(c=-3\)，则 \(a>b>c\)，但 \(a+b=-3=c\)，不满足 \(a+b>c\). 因此这组整数可以说明原命题是假命题.
\end{solution}

\begin{problem}
某学习小组由学生和教师组成，人员构成同时满足以下三个条件：

\begin{circlelist}
\item 男学生人数多于女学生人数；
\item 女学生人数多于教师人数；
\item 教师人数的两倍多于男学生人数.
\end{circlelist}
\begin{enumerate}
\item 若教师人数为 \(4\)，则女学生人数的最大值为\fillinblank{}.
\item 该小组人数的最小值为\fillinblank{}.
\end{enumerate}
\end{problem}

\begin{answer}
6；12
\end{answer}

\begin{solution}
设男学生、女学生和教师人数分别为 \(m\)，\(f\)，\(t\)，则 \(m>f>t\)，且 \(m<2t\)，并且 \(m\)，\(f\)，\(t\in\Z_{\ge0}\).

\begin{enumerate}
\item 当 \(t=4\) 时，\(m<8\)，所以 \(m\le7\). 又 \(f<m\)，故 \(f\le6\). 取 \(m=7\)，\(f=6\) 时三个条件均满足，所以女学生人数的最大值为 \(6\).
\item 由人数非负且 \(f>t\)，若 \(t=0\)，则 \(f\ge1\)、\(m\ge2\)，但 \(m<2t=0\)，矛盾，故 \(t\ge1\). 由 \(m>f>t\) 得 \(m\ge t+2\)，而 \(m<2t\). 当 \(t=1\) 或 \(t=2\) 时分别得到 \(m\ge3\) 且 \(m<2\)，或 \(m\ge4\) 且 \(m<4\)，均不可能，因此 \(t\ge3\). 从而 \(m+f+t\ge(t+2)+(t+1)+t\ge12\). 取 \(t=3\)，\(f=4\)，\(m=5\) 时条件成立，故该小组人数的最小值为 \(12\).
\end{enumerate}
\end{solution}

\section{解答题}

\begin{problem}

已知等差数列 \(\{a_n\}\) 和等比数列 \(\{b_n\}\) 满足 \(a_1=b_1=1\)，\(a_2+a_4=10\)，\(b_2b_4=a_5\).

\begin{enumerate}
\item 求 \(\{a_n\}\) 的通项公式；
\item 求和： \(b_1+b_3+b_5+\cdots+b_{2n-1}\).
\end{enumerate}
\end{problem}

\begin{solution}
设等差数列 \(\{a_n\}\) 的公差为 \(d\)，则
\(a_2+a_4=(1+d)+(1+3d)=10\)，解得 \(d=2\). 因此
\(a_n=a_1+(n-1)d=1+2(n-1)=2n-1\).

设等比数列 \(\{b_n\}\) 的公比为 \(q\). 由上式得 \(a_5=1+4\times2=9\)，而 \(b_1=1\)，所以
\(b_2b_4=q\cdot q^3=q^4=9\)，从而 \(q^2=3\). 不论 \(q=\sqrt3\) 还是 \(q=-\sqrt3\)，都有
\(b_{2k-1}=q^{2k-2}=(q^2)^{k-1}=3^{k-1}\).
因此
\(b_1+b_3+b_5+\cdots+b_{2n-1}=\sum_{k=1}^{n}3^{k-1}=\dfrac{3^n-1}{3-1}=\dfrac{3^n-1}{2}\).
\end{solution}

\begin{problem}

已知函数 \(f(x)=\sqrt3\cos\left(2x-\dfrac\pi3\right)-2\sin x\cos x\).

\begin{enumerate}
\item 求 \(f(x)\) 的最小正周期；
\item 求证：当 \(x\in\left[-\dfrac\pi4,\dfrac\pi4\right]\) 时，\(f(x)\ge-\dfrac12\).
\end{enumerate}
\end{problem}

\begin{solution}
利用和角公式及二倍角公式，得
\(\sqrt3\cos\left(2x-\dfrac\pi3\right)=\dfrac{\sqrt3}{2}\cos2x+\dfrac32\sin2x\)，且 \(2\sin x\cos x=\sin2x\). 所以
\(f(x)=\dfrac{\sqrt3}{2}\cos2x+\dfrac12\sin2x=\sin\left(2x+\dfrac\pi3\right)\).

\begin{enumerate}
\item 因为 \(\sin(2x+\dfrac\pi3)\) 的角频率为 \(2\)，所以其最小正周期为 \(T=\dfrac{2\pi}{2}=\pi\).
\item 当 \(x\in\left[-\dfrac\pi4,\dfrac\pi4\right]\) 时，\(2x+\dfrac\pi3\in\left[-\dfrac\pi6,\dfrac{5\pi}6\right]\). 正弦函数在该区间上的最小值为 \(\sin(-\dfrac\pi6)=-\dfrac12\)，故 \(f(x)=\sin(2x+\dfrac\pi3)\ge-\dfrac12\).
\end{enumerate}
\end{solution}

\begin{problem}

某大学艺术专业 \(400\) 名学生参加某次测评，根据男女学生人数比例，使用分层抽样的方法从中随机抽取了 \(100\) 名学生，记录他们的分数，将数据分成 \(7\) 组：\(\left[20,30\right)\)，\(\left[30,40\right)\)，\(\cdots\)，\(\left[80,90\right]\)，并整理得到如下频率分布直方图：

\begin{enumerate}
\item 从总体的 \(400\) 名学生中随机抽取一人，估计其分数小于 \(70\) 的概率；
\item 已知样本中分数小于 \(40\) 的学生有 \(5\) 人，试估计总体中分数在区间 \(\left[40,50\right)\) 内的人数；
\item 已知样本中有一半男生的分数不小于 \(70\)，且样本中分数不小于 \(70\) 的男女生人数相等. 试估计总体中男生和女生人数的比例.
\end{enumerate}

\texfigure[max width=0.62\linewidth]{img_repaint/2017/beijing_liberal/q17_fig1.tex}
\end{problem}

\begin{solution}
\begin{enumerate}
\item 由直方图知，分数不小于 \(70\) 的样本频率为 \((0.04+0.02)\times10=0.6\). 因此分数小于 \(70\) 的频率为 \(1-0.6=0.4\)，据此估计从总体中随机抽取一名学生，其分数小于 \(70\) 的概率为 \(0.4\).
\item 设区间 \([40,50)\) 对应的频率为 \(p\). 样本中分数小于 \(40\) 的学生有 \(5\) 人，所以这部分频率为 \(\dfrac5{100}=0.05\). 由直方图，其他各组的频率为 \(0.01\times10\)，\(0.02\times10\)，\(0.04\times10\)，\(0.02\times10\)，故
\(0.05+10p+0.01\times10+0.02\times10+0.04\times10+0.02\times10=1\).
解得 \(p=0.05\). 因此估计总体中分数在区间 \([40,50)\) 内的人数为 \(400\times0.05=20\) 人.
\item 样本中分数不小于 \(70\) 的频率为 \(0.6\). 由于其中男生和女生人数相等，男生中分数不小于 \(70\) 的频率为 \(0.3\). 又这部分男生占全体男生的一半，所以男生在样本中的频率为 \(0.3\div\dfrac12=0.6\)，女生在样本中的频率为 \(1-0.6=0.4\). 分层抽样保持总体中的男女比例，因此总体中男生和女生人数的比例约为 \(0.6:0.4=3:2\).
\end{enumerate}
\end{solution}

\begin{problem}

如图，在三棱锥 \(P-ABC\) 中，\(PA\perp AB\)，\(PA\perp BC\)，\(AB\perp BC\)，\(PA=AB=BC=2\)，\(D\) 为线段 \(AC\) 的中点，\(E\) 为线段 \(PC\) 上一点.

\begin{enumerate}
\item 求证：\(PA\perp BD\)；
\item 求证：平面 \(BDE\perp\) 平面 \(PAC\)；
\item 当 \(PA\parallel\) 平面 \(BDE\) 时，求三棱锥 \(E-BCD\) 的体积.
\end{enumerate}

\texfigure[max width=0.58\linewidth]{img_repaint/2017/beijing_liberal/q18_fig1.tex}
\end{problem}

\begin{solution}
\begin{enumerate}
\item 因为 \(PA\perp AB\)，\(PA\perp BC\)，且直线 \(AB\)，\(BC\) 是平面 \(ABC\) 内相交的两条直线，所以 \(PA\perp\) 平面 \(ABC\). 又 \(BD\subset\) 平面 \(ABC\)，故 \(PA\perp BD\).
\item 在等腰三角形 \(ABC\) 中，\(AB=BC\)，且 \(D\) 是底边 \(AC\) 的中点，所以中线 \(BD\perp AC\). 由第（1）问知 \(BD\perp PA\). 直线 \(PA\)，\(AC\) 是平面 \(PAC\) 内相交的两条直线，因此 \(BD\perp\) 平面 \(PAC\). 又 \(BD\subset\) 平面 \(BDE\)，所以平面 \(BDE\perp\) 平面 \(PAC\).
\item 因为 \(PA\parallel\) 平面 \(BDE\)，且 \(PA\subset\) 平面 \(PAC\)，而两平面 \(PAC\) 与 \(BDE\) 的交线为 \(DE\)，所以 \(PA\parallel DE\). 在三角形 \(PAC\) 中，\(D\) 是 \(AC\) 的中点，过 \(D\) 的直线 \(DE\) 平行于 \(PA\)，故 \(E\) 是 \(PC\) 的中点，且 \(DE=\dfrac12PA=1\).
又因为 \(PA\perp\) 平面 \(ABC\)，所以 \(DE\perp\) 平面 \(ABC\). 由 \(AB\perp BC\) 及 \(AB=BC=2\)，得 \(S_{\triangle ABC}=\dfrac12\times2\times2=2\). 点 \(D\) 是 \(AC\) 的中点，故 \(S_{\triangle BCD}=\dfrac12S_{\triangle ABC}=1\). 于是
\(V_{E-BCD}=\dfrac13S_{\triangle BCD}\cdot DE=\dfrac13\times1\times1=\dfrac13\).
\end{enumerate}
\end{solution}

\begin{problem}

已知椭圆 \(C\) 的两个顶点分别为 \(A(-2,0)\)，\(B(2,0)\)，焦点在 \(x\) 轴上，离心率为 \(\dfrac{\sqrt3}{2}\).

\begin{enumerate}
\item 求椭圆 \(C\) 的方程；
\item 点 \(D\) 为 \(x\) 轴上一点，过 \(D\) 作 \(x\) 轴的垂线交椭圆 \(C\) 于不同的两点 \(M\)，\(N\)，过 \(D\) 作 \(AM\) 的垂线交 \(BN\) 于点 \(E\). 求证：\(\triangle BDE\) 与 \(\triangle BDN\) 的面积之比为 \(4:5\).
\end{enumerate}
\end{problem}

\begin{solution}
\begin{enumerate}
\item 由椭圆的顶点为 \(A(-2,0)\)，\(B(2,0)\)，且焦点在 \(x\) 轴上，得长半轴 \(a=2\). 由离心率 \(e=\dfrac ca=\dfrac{\sqrt3}{2}\)，得焦半距 \(c=\sqrt3\). 因此短半轴的平方为 \(b^2=a^2-c^2=4-3=1\)，所以椭圆 \(C\) 的方程为 \(\dfrac{x^2}{4}+y^2=1\).
\item 设 \(D=(x_0,0)\)，由直线与椭圆交于不同两点，得 \(-2<x_0<2\). 不妨设 \(M=(x_0,y_0)\)，\(N=(x_0,-y_0)\)，其中 \(y_0>0\). 因为 \(M\)，\(N\) 在椭圆上，所以
\(\dfrac{x_0^2}{4}+y_0^2=1\)，即 \(4-x_0^2=4y_0^2\).

令 \(A_0=\dfrac{x_0+2}{y_0}\)，\(B_0=\dfrac{y_0}{2-x_0}\). 则直线 \(DE\) 的斜率为 \(-A_0\)，直线 \(BN\) 的斜率为 \(B_0\)，两直线方程分别为 \(y=-A_0(x-x_0)\) 和 \(y=B_0(x-2)\). 设 \(E=(x_E,y_E)\)，联立两式得
\(\lvert y_E\rvert=\dfrac{A_0B_0(2-x_0)}{A_0+B_0}=\dfrac{A_0y_0}{A_0+B_0}\).
又
\(A_0+B_0=\dfrac{x_0+2}{y_0}+\dfrac{y_0}{2-x_0}=\dfrac{(x_0+2)(2-x_0)+y_0^2}{y_0(2-x_0)}=\dfrac{5y_0}{2-x_0}\).
所以
\(\dfrac{\lvert y_E\rvert}{y_0}=\dfrac{A_0}{A_0+B_0}=\dfrac{(x_0+2)(2-x_0)}{5y_0^2}=\dfrac{4-x_0^2}{5y_0^2}=\dfrac45\).
三角形 \(BDE\) 与 \(BDN\) 具有同一条底边 \(BD\)，故面积之比等于对应高之比，即
\(S_{\triangle BDE}:S_{\triangle BDN}=\lvert y_E\rvert:y_0=4:5\).
\end{enumerate}
\end{solution}

\begin{problem}

已知函数 \(f(x)=\e^x\cos x-x\).

\begin{enumerate}
\item 求曲线 \(y=f(x)\) 在点 \((0,f(0))\) 处的切线方程；
\item 求函数 \(f(x)\) 在区间 \(\left[0,\dfrac\pi2\right]\) 上的最大值和最小值.
\end{enumerate}
\end{problem}

\begin{solution}
\begin{enumerate}
\item \(f'(x)=\e^x(\cos x-\sin x)-1\). 因为 \(f(0)=1\)，且 \(f'(0)=1\times(1-0)-1=0\)，所以曲线在点 \((0,f(0))=(0,1)\) 处的切线斜率为 \(0\)，切线方程为 \(y=1\).
\item 令 \(g(x)=f'(x)=\e^x(\cos x-\sin x)-1\)，则
\(g'(x)=\e^x(\cos x-\sin x)+\e^x(-\sin x-\cos x)=-2\e^x\sin x\).
当 \(x\in\left[0,\dfrac\pi2\right]\) 时，\(g'(x)\le0\)，所以 \(g\) 在该区间上单调递减. 又 \(g(0)=0\)，故对区间内的 \(x\)，有 \(f'(x)=g(x)\le0\)，函数 \(f\) 在该区间上单调递减.
因此最大值为 \(f(0)=1\)，最小值为 \(f\left(\dfrac\pi2\right)=\e^{\frac\pi2}\cos\dfrac\pi2-\dfrac\pi2=-\dfrac\pi2\).
\end{enumerate}
\end{solution}
